\chapter{Entire Series, Coleman's Resultant, and the Riesz--Coleman Decomposition}
\label{chap:riesz-coleman}

Chapter~\ref{chap:compact} attached to a compactoid operator $u$ its Fredholm determinant
$\det(1-Tu)$ and proved Serre's theorems: $1-au$ is invertible exactly when $\det(1-Tu)$ is a
unit at $a$ (Theorem~\ref{thm:serre11}), a zero of the determinant is a reciprocal eigenvalue
(Theorem~\ref{thm:serre12}), and each zero carries a finite-dimensional Riesz summand
(Theorem~\ref{thm:riesz}).  Serre's decomposition is attached to a \emph{single} zero, and it
lives over a field.

The applications need more.  A slope decomposition splits off \emph{all} the eigenvalues of a
given valuation at once, which means splitting the determinant as $\det(1-Tu) = Q\cdot S$ with
$Q$ a polynomial and asking for the corresponding splitting of the space; and the coefficient
ring in the halo application of Chapter~\ref{chap:lwx-halo} is the Banach--Tate ring
$\Lhalo[1/T]$, not a field, and is not Noetherian.  This chapter formalises the general
statement over a Banach--Tate ring, following Johansson--Newton \cite{JN} Theorem~2.2.2, which
is Buzzard \cite{Buzzard} Theorem~3.3, which is Bella\"iche \cite{Bellaiche} Theorem~II.2.18,
and ultimately Coleman \cite{Coleman} Theorem~A4.3.  No Noetherian hypothesis appears anywhere:
none of the sources' proofs uses one.

The route is Bella\"iche's.  Given the factorisation $\det(1-Tu) = QS$, one replaces $u$ by the
auxiliary operator $\varphi' = 1 - Q^{*}(u)/Q^{*}(0)$ and computes \emph{its} Fredholm
determinant by a spectral mapping formula; that determinant turns out to have a zero of order
$\deg Q$ at $1$, and Serre's projector at that zero is the projector one wants.  So we need
three things first: the ring $R\{\{T\}\}$ of entire series and division inside it (\S1),
Coleman's resultant $D(B,P)$ and the spectral mapping formula (\S2), and then the decomposition
itself (\S3).  Section~4 adds the vertex factorisation that produces the polynomial $Q$ in the
first place.

\section{Entire series and division}

\begin{definition}
  \label{def:isentire}
  \uses{def:isrestricted}
  \lean{PowerSeries.IsEntire, PowerSeries.entireSubring}
  \leanok
  A power series $F = \sum a_n T^n$ over a normed ring $R$ is \emph{entire} if
  $\lVert a_n\rVert C^n \to 0$ for every $C > 0$; equivalently, if it is restricted at every
  radius in the sense of Definition~\ref{def:isrestricted}.  The entire series form a subring
  $R\{\{T\}\}$ of $R\llbracket T\rrbracket$.  A \emph{Fredholm series} is an entire series with
  constant term $1$.
\end{definition}

This is \cite[Definition II.1.16]{Bellaiche} and the ring written $R\{\{T\}\}$ in
\cite[\S2.1]{JN}; Theorem~\ref{thm:charpowerseries-entire} says the Fredholm determinant of a
compactoid operator lies in it.

\begin{theorem}[Euclidean division in $R\{\{T\}\}$, {\cite[Proposition II.2.8]{Bellaiche}}]
  \label{thm:entire-division}
  \uses{def:isentire, thm:wdivision, thm:wdivision-unique}
  \lean{PowerSeries.IsEntire.exists_eq_mul_add, PowerSeries.eq_mul_add_q_unique,
        PowerSeries.eq_mul_add_r_unique}
  \leanok
  Let $g \in R[T]$ have invertible leading coefficient and degree $n$.  Every entire $F$ can be
  written uniquely as $F = g\,q + r$ with $q$ entire and $r$ a polynomial of degree $< n$.
\end{theorem}

\begin{proof}
  \leanok
  \uses{thm:wdivision, thm:wdivision-unique}
  Existence and uniqueness are Martin's Weierstrass division (Theorem~\ref{thm:wdivision} and
  Theorem~\ref{thm:wdivision-unique}) applied at a radius large enough that $g$ dominates: the
  divisor of a distinguished division is precisely a polynomial whose leading coefficient
  achieves the Gauss norm, and for a polynomial with invertible leading coefficient this happens
  at all sufficiently large radii.  Entireness of the quotient is then the statement that the
  division is performed at every radius simultaneously, which uniqueness supplies.
\end{proof}

\begin{definition}
  \label{def:isentirecoprime}
  \uses{def:isentire}
  \lean{PowerSeries.IsEntireCoprime}
  \leanok
  Two entire series are \emph{relatively prime in $R\{\{T\}\}$} if the ideal they generate is
  the unit ideal (\cite[Definition 2.2.1]{JN}).
\end{definition}

\begin{lemma}[{\cite[Corollary II.2.9]{Bellaiche}}]
  \label{lem:coprime-remainder}
  \uses{def:isentirecoprime, thm:entire-division}
  \lean{PowerSeries.isEntireCoprime_iff_isCoprime_of_eq_mul_add}
  \leanok
  If $F = g\,q + r$ is the division of Theorem~\ref{thm:entire-division}, then $F$ and $g$ are
  relatively prime in $R\{\{T\}\}$ if and only if $r$ and $g$ are relatively prime in $R[T]$.
\end{lemma}

Division reduces coprimality questions to polynomial ones; that is the only use we make of it.

\begin{definition}
  \label{def:isgoodzero}
  \uses{def:isentire}
  \lean{PowerSeries.IsGoodZero, PowerSeries.hasseDeriv_mul}
  \leanok
  An entire $F$ has a \emph{good zero of order $s$} at $a$ if the divided derivatives
  $\Delta^{i}F$ vanish at $a$ for $i < s$ and $\Delta^{s}F(a)$ is a unit
  (\cite[Definition II.2.11]{Bellaiche}).  The divided derivatives obey the Leibniz rule
  $\Delta^{s}(FG) = \sum_{i+j=s}\Delta^{i}F\cdot\Delta^{j}G$, so good zeros multiply.
\end{definition}

\section{Coleman's $D(B,P)$ and the spectral mapping formula}

Coleman's construction \cite[\S A3--A4]{Coleman} takes a polynomial $B$ and a Fredholm series
$P$ and produces the series whose reciprocal roots are the $B$-images of those of $P$:
informally $D(B,P)(T) = \prod_i (1 - T\,B(t_i))$ over the roots $t_i$ of the reciprocal
polynomial $P^{*}$.  Written that way the definition needs roots; written as a resultant it
needs nothing.

\begin{definition}
  \label{def:dpoly}
  \lean{TateFredholm.Coleman.dPoly, TateFredholm.Coleman.bQ}
  \leanok
  For polynomials $B, P$ over a commutative ring and $N \geq \deg P$, set
  \[
    D_N(B,P) \;=\; \Res_X\bigl(P^{*}, \; 1 - T\,B(X)\bigr) \in R[T],
  \]
  the resultant of the reversed polynomial $P^{*} = \mathrm{reflect}_N P$ against
  $1 - T B(X)$.  We also write $b_Q = 1 - u^{-1}Q^{*}$ for a polynomial $Q$ with leading
  coefficient the unit $u$; this is Bella\"iche's auxiliary polynomial.
\end{definition}

\begin{lemma}[{\cite[Lemma II.2.13]{Bellaiche}}]
  \label{lem:dpoly-mul}
  \uses{def:dpoly}
  \lean{TateFredholm.Coleman.dPoly_mul}
  \leanok
  $D_N(B, PQ) = D_N(B,P)\,D_N(B,Q)$: the construction is multiplicative in its second argument.
\end{lemma}

\begin{proof}
  \leanok
  Multiplicativity of the resultant in its first argument, together with
  $\mathrm{reflect}(PQ) = \mathrm{reflect}(P)\,\mathrm{reflect}(Q)$ at compatible degrees.
\end{proof}

\begin{proposition}[The finite-dimensional case]
  \label{prop:resultant-charpoly}
  \uses{def:dpoly}
  \lean{Matrix.resultant_charpoly, Matrix.det_aeval_eq_prod_roots}
  \leanok
  For a square matrix $A$ over any commutative ring and any polynomial $g$,
  \[
    \Res(\mathrm{charpoly}\,A,\; g) \;=\; \det\bigl(g(A)\bigr).
  \]
\end{proposition}

\begin{proof}
  \leanok
  Over an algebraically closed field both sides equal $\prod_\lambda g(\lambda)$ over the
  eigenvalues with multiplicity.  The general case follows from the universal characteristic
  polynomial: both sides are polynomial identities in the matrix entries and the coefficients of
  $g$, so it suffices to verify them over $\Z[x_{ij}, y_k]$, which embeds in an algebraic closure
  of its fraction field.
\end{proof}

\begin{definition}
  \label{def:dseries}
  \uses{def:dpoly, def:isentire}
  \lean{TateFredholm.Coleman.dSeries, TateFredholm.Coleman.norm_coeff_dSeries_sub_le,
        PowerSeries.IsEntire.dSeries}
  \leanok
  For a Fredholm series $F$ over a Banach--Tate ring, $D(B,F)$ is the coefficientwise limit of
  $D_N(B, \mathrm{trunc}_N F)$.  The limit exists because each coefficient of $D_N$ is a fixed
  polynomial in the coefficients of the truncation, and that polynomial is Lipschitz in them;
  $D(B, \cdot)$ is multiplicative and carries entire series to entire series.
\end{definition}

Entireness uses Buzzard's rescaling trick \cite[p.~21]{Buzzard}: the identity
$\Res(Q,P) = \Res(u^{-n}Q(uT), P(uT))$ lets one renormalise the radius at which the estimate is
performed.

\begin{lemma}[{\cite[Lemma II.2.14 and Proposition II.2.15]{Bellaiche}}]
  \label{lem:dseries-goodzero}
  \uses{def:dseries, def:isentirecoprime, def:isgoodzero}
  \lean{TateFredholm.Coleman.isUnit_evalT_one_dSeries_bQ_iff,
        TateFredholm.Coleman.isGoodZero_dSeries_bQ}
  \leanok
  Let $Q$ be a polynomial with unit leading coefficient and constant term $1$, and $S$ a
  Fredholm series.  Then $D(b_Q, S)(1)$ is a unit if and only if $Q$ and $S$ are relatively
  prime in $R\{\{T\}\}$; and $D(b_Q, QS)$ has a good zero of order $\deg Q$ at $1$.
\end{lemma}

\begin{proof}
  \leanok
  \uses{lem:dpoly-mul, lem:coprime-remainder, def:isgoodzero}
  By Lemma~\ref{lem:dpoly-mul}, $D(b_Q, QS) = D(b_Q, Q)\,D(b_Q, S)$.  The first factor is
  computed directly: $b_Q$ was built so that $D(b_Q, Q)$ is $(1-T)^{\deg Q}$ up to a unit, which
  is a good zero of order $\deg Q$ at $1$.  The second factor is a unit at $1$ by the
  coprimality criterion, which is Lemma~\ref{lem:coprime-remainder} transported through the
  resultant.  Good zeros multiply by the Leibniz rule of Definition~\ref{def:isgoodzero}.
\end{proof}

\begin{theorem}[Spectral mapping, {\cite[Proposition II.2.16]{Bellaiche}}]
  \label{thm:spectral-mapping}
  \uses{def:dseries, prop:resultant-charpoly, def:iscompactoid}
  \lean{TateFredholm.Coleman.charPowerSeries_aeval}
  \leanok
  For a compactoid $\varphi$ on $c(I,R)$ and any polynomial $B$ with $B(0)=0$,
  \[
    \det\bigl(1 - T\,B(\varphi)\bigr) \;=\; D\bigl(B, \det(1-T\varphi)\bigr).
  \]
\end{theorem}

\begin{proof}
  \leanok
  \uses{prop:resultant-charpoly, def:dseries}
  For a finite matrix this is Proposition~\ref{prop:resultant-charpoly} rewritten in reversed
  form.  A compactoid operator is a limit of its truncations to finite blocks, both sides are
  continuous in the operator coefficientwise, and $D(B,\cdot)$ was defined as exactly that
  limit.
\end{proof}

\section{The Riesz--Coleman decomposition}

\begin{theorem}[{\cite[Theorem 2.2.2]{JN}; \cite[Theorem 3.3]{Buzzard};
                 \cite[Theorem A4.3]{Coleman}}]
  \label{thm:riesz-coleman}
  \uses{thm:spectral-mapping, lem:dseries-goodzero, thm:riesz, def:iscompactoid}
  \lean{TateFredholm.exists_rieszColemanProjection, TateFredholm.IsRieszColemanProjection}
  \leanok
  Let $R$ be a Banach--Tate ring, $u$ compactoid on $M = c(I,R)$, and suppose
  $\det(1-Tu) = Q\,S$ with $Q$ a polynomial with unit leading coefficient and constant term $1$,
  relatively prime to the Fredholm series $S$ in $R\{\{T\}\}$.  Then there is a continuous
  projector $p$, lying in the closure of $R[u]$, with $Q^{*}(u)\,w = p$ for some $w$ in that
  closure and $Q^{*}(u)^{\deg Q}(1-p) = 0$.
\end{theorem}

\begin{proof}
  \leanok
  \uses{thm:spectral-mapping, lem:dseries-goodzero, thm:riesz}
  Put $\varphi' = 1 - Q^{*}(u)/Q^{*}(0)$, a compactoid operator.  Theorem~\ref{thm:spectral-mapping}
  computes $\det(1 - T\varphi') = D(b_Q, \det(1-Tu))$, and Lemma~\ref{lem:dseries-goodzero}
  says that this has a good zero of order $\deg Q$ at $1$.  Serre's projector at that zero
  (the construction behind Theorem~\ref{thm:riesz}) is $p$.  Because the projector is produced
  as an operator-norm limit of polynomials in $\varphi'$, and $\varphi'$ is itself a polynomial
  in $u$, both $p$ and the inverse witness $w$ lie in the closure of $R[u]$ --- so they commute
  with $u$, not merely with $\varphi'$.
\end{proof}

The refinements of \cite[Proposition 3.2 and Theorem 3.3]{Buzzard} and the last sentence of
\cite[Theorem 2.2.2]{JN} follow.  We record them as the API of the predicate
$\mathrm{IsRieszColemanProjection}$.

\begin{theorem}[The full decomposition]
  \label{thm:riesz-coleman-full}
  \uses{thm:riesz-coleman}
  \lean{TateFredholm.IsRieszColemanProjection.finite,
        TateFredholm.IsRieszColemanProjection.projective,
        TateFredholm.IsRieszColemanProjection.rankAtStalk,
        TateFredholm.IsRieszColemanProjection.charPowerSeries_mul_one_sub,
        TateFredholm.IsRieszColemanProjection.charPowerSeries_mul,
        TateFredholm.IsRieszColemanProjection.ker_aeval_reverse,
        TateFredholm.IsRieszColemanProjection.isUnit_mul_one_sub_add,
        TateFredholm.IsRieszColemanProjection.eq_range_of_isTopCompl}
  \leanok
  With the notation of Theorem~\ref{thm:riesz-coleman}, write $N = \ker Q^{*}(u)$.  Then
  $M = N \oplus \mathrm{im}\,p$ topologically; $N$ is finitely generated projective of rank
  $\deg Q$ at every prime; $u$ is invertible on $N$; $\det(1-Tu \mid N) = Q$ and
  $\det(1-Tu\mid \mathrm{im}\,p) = S$; and $\mathrm{im}\,p$ is the unique $u$-stable closed
  complement on which $Q^{*}(u)$ is invertible.
\end{theorem}

\begin{lemma}[{\cite[Lemma 3.1]{Buzzard}}]
  \label{lem:coprime-iff-unit}
  \uses{def:isentirecoprime}
  \lean{TateFredholm.isEntireCoprime_iff_isUnit_aeval_reverse}
  \leanok
  $Q$ and $S$ are relatively prime in $R\{\{T\}\}$ if and only if $Q^{*}(u)$ is invertible on
  the summand where the determinant is $S$.
\end{lemma}

\begin{proposition}[Determinants of summands]
  \label{prop:det-summand}
  \uses{thm:riesz-coleman-full}
  \lean{TateFredholm.charPowerSeries_eq_mul_of_comm,
        TateFredholm.exists_polynomial_charPowerSeries_of_range_le,
        TateFredholm.exists_matrix_realisation}
  \leanok
  If $u$ commutes with a continuous projector then its determinant is the product of the
  determinants of the two pieces; on a finitely generated summand the determinant is a
  polynomial, computed as the reversed characteristic polynomial of any matrix realisation.
\end{proposition}

The matrix-level facts behind Proposition~\ref{prop:det-summand} are elementary but were not in
mathlib.

\begin{lemma}
  \label{lem:charpolyrev}
  \lean{Matrix.charpolyRev_map, Matrix.charpolyRev_mul_comm, Matrix.charpolyRev_eq_of_mul_eq,
        Matrix.exists_mul_eq_of_idempotent, Matrix.charpolyRev_eq_one_sub_pow_rank}
  \leanok
  Write $\mathrm{charpolyRev}\,M = \det(1 - XM)$.  It commutes with ring homomorphisms and
  satisfies Sylvester's identity $\det(1-XAB) = \det(1-XBA)$; an idempotent matrix with free
  range factors as $E = PQ$, $QP = 1$, and along such a factorisation
  $\mathrm{charpolyRev}\,M = \mathrm{charpolyRev}(QMP)$ whenever $ME = EM = M$.  Over a field,
  if moreover $M = E - N$ with $N$ nilpotent then
  $\mathrm{charpolyRev}\,M = (1-X)^{\mathrm{rank}\,E}$.
\end{lemma}

\section{The vertex factorisation}

Theorem~\ref{thm:riesz-coleman} takes the factorisation $\det(1-Tu) = QS$ as input.  Producing
it is a Newton-polygon statement: a vertex of the polygon at which the coefficient is a unit
splits off a polynomial factor of the corresponding degree.

\begin{definition}
  \label{def:dominant-index}
  \uses{def:isentire}
  \lean{PowerSeries.IsDominantIndex, PowerSeries.IsDominantPoly}
  \leanok
  The index $N$ is \emph{dominant} for $F = \sum a_k T^k$ at radius $\rho$ if
  $\lVert a_k\rVert\rho^k \leq \lVert a_N\rVert\rho^N$ for all $k$, with strict inequality for
  $k > N$; that is, $N$ is Bella\"iche's largest integer minimising $v(a_n) - n\nu$
  (\cite[Definition II.2.3]{Bellaiche}).  A polynomial of degree $N$ dominant at $\rho$ with
  invertible leading coefficient is a \emph{$\rho$-dominant polynomial}
  (\cite[Definition II.2.4]{Bellaiche}).
\end{definition}

\begin{lemma}
  \label{lem:dominant-distinguished}
  \uses{def:dominant-index, def:ismuldistinguished}
  \lean{PowerSeries.IsDominantPoly.isMulDistinguished, PowerSeries.IsDominantPoly.mul,
        PowerSeries.exists_isDominantPoly_of_isUnit_leadingCoeff}
  \leanok
  A $\rho$-dominant polynomial is exactly a Martin-distinguished element at radius $\rho$
  (Definition~\ref{def:ismuldistinguished}); dominant polynomials are closed under
  multiplication.
\end{lemma}

That identification is what lets Chapter~\ref{chap:wp} be reused here.

\begin{theorem}[The slope factorisation, {\cite[Theorem II.3.6]{Bellaiche}};
                {\cite[Definition 2.2.5]{JN}}]
  \label{thm:slope-factorisation}
  \uses{def:dominant-index, lem:dominant-distinguished, thm:wdivision, def:isentirecoprime}
  \lean{PowerSeries.exists_isDominantFactorization, PowerSeries.IsDominantFactorization,
        PowerSeries.IsDominantFactorization.isEntireCoprime}
  \leanok
  Let $F$ be entire over a Banach--Tate ring, with dominant index $N$ at radius $\rho$ and with
  $\lVert\cdot\rVert$ multiplicative on the dominant coefficient $a_N$.  Then $F = PG$ with $P$
  a $\rho$-dominant polynomial of degree $N$ and $G$ entire of dominant index $0$ at $\rho$;
  moreover $P$ and $G$ are relatively prime in $R\{\{T\}\}$.
\end{theorem}

\begin{proof}
  \leanok
  \uses{lem:dominant-distinguished, thm:wdivision, thm:wprep}
  Bella\"iche and Kedlaya \cite[Proposition 3.2.2]{Kedlaya} both run a Newton iteration
  $P_{l+1} = P_l + X_l$, $Q_{l+1} = Q_l + Y_l$ driven by the division estimates
  $\lVert f\rVert = \max(\lVert g\rVert\lVert q\rVert, \lVert r\rVert)$.  We instead observe that
  those division estimates \emph{are} Martin's, so the hypotheses say exactly that $F$ is
  distinguished of order $N$ at radius $\rho$ (Lemma~\ref{lem:dominant-distinguished}); dividing
  $T^N$ by $F$ (Theorem~\ref{thm:wdivision}) and subtracting the remainder from $T^N$ produces
  the monic factor.  This is Weierstrass preparation (Theorem~\ref{thm:wprep}) in the form
  needed.  For coprimality (the norm-level form of \cite[Lemma 2.2.7]{JN}, and where \cite{JN}
  argue with resultants over the Gelfand spectrum): $\lVert G-1\rVert_\rho < 1$, so $G$ is a unit
  of the Tate algebra of radius $\rho$; dividing $G^{-1}$ by $P$ there gives
  $1 = Pq' + rG$ with $\deg r < N$, and the cofactor is entire because the entire division of
  $1 - rG$ by $P$ has the same quotient, by uniqueness.  That is a B\'ezout identity for $P$ and
  $G$ in $R\{\{T\}\}$.
\end{proof}

Together with Theorem~\ref{thm:riesz-coleman-full}, Theorem~\ref{thm:slope-factorisation} says:
\emph{a vertex of the Newton polygon of $\det(1-Tu)$ whose coefficient is a multiplicative unit
splits off a $u$-stable summand of rank equal to the abscissa of the vertex, carrying exactly
the eigenvalues of smaller slope.}  This is the form in which the halo $U_p$ of
\S\ref{sec:tate-riesz} is decomposed.

\section{Auxiliary constructions}

Three further pieces of the model-space toolkit are used in
Chapters~\ref{chap:lwx-halo}--\ref{chap:lwx-slopes} and recorded here.

\begin{lemma}[Two-sided weight bounds]
  \label{lem:two-sided}
  \uses{lem:hadamard, def:minor}
  \lean{TateFredholm.norm_charCoeff_le_pow_two_sided}
  \leanok
  Suppose the matrix of $u$ satisfies $\lVert u_{ab}\rVert \leq \sigma^{\max(w(a) - w'(b),\,0)}$
  for weights $w, w'$.  Then $\lVert c_n(u)\rVert \leq \sigma^{f(n)}$ whenever
  $f(n) \leq \sum_{a\in S} w(a) - \sum_{a \in S} w'(a)$ for all $\lvert S\rvert = n$.
\end{lemma}

\begin{proof}
  \leanok
  \uses{lem:hadamard}
  A permutation $\pi$ of $S$ satisfies $\sum_{a\in S} w'(\pi a) = \sum_{a \in S} w'(a)$, so every
  Leibniz monomial of $\mathrm{minor}(u,S)$ carries the same column contribution and the row
  contribution is $\sum_{a\in S}w(a)$.  This is the Hadamard bound of Lemma~\ref{lem:hadamard}
  with the column weight subtracted; note that nothing is inverted, so the estimate runs over an
  arbitrary complete ultrametric ring rather than a field.
\end{proof}

This is the estimate behind \cite[Theorem 3.16]{LWX}, where $\sigma = \lVert T\rVert$ and the
conjugation by $\mathrm{diag}(T^{w})$ that \cite{LWX} performs explicitly is replaced by the
observation that principal minors are invariant under diagonal conjugation.

\begin{lemma}[Diagonal intertwining]
  \label{lem:diag-intertwine}
  \uses{def:minor}
  \lean{TateFredholm.minor_eq_of_diag_intertwine, TateFredholm.charCoeff_eq_of_diag_intertwine,
        TateFredholm.charPowerSeries_eq_of_diag_intertwine}
  \leanok
  If $D\,M_v = M_u\,D$ entrywise for a diagonal $D$ with unit entries then
  $\mathrm{minor}(v,S) = \mathrm{minor}(u,S)$ for every finite $S$, and hence
  $\det(1-Tv) = \det(1-Tu)$.  No compactness hypothesis and no inversion of $D$ is needed.
\end{lemma}

The point of avoiding $D^{-1}$ is that the diagonal occurring in \cite[Proposition 2.17]{LWX} is
$\mathrm{diag}(n!)$, whose inverse is unbounded.

\begin{definition}[Block-diagonal and rectangular block maps]
  \label{def:blockmaps}
  \uses{prop:blockop}
  \lean{TateFredholm.blockDiag, TateFredholm.diagBlockEquiv, TateFredholm.blockOpMap,
        TateFredholm.blockMap, TateFredholm.blockMapEquiv}
  \leanok
  An operator $f$ on $c(I,R)$ induces the block-diagonal operator $\mathrm{blockDiag}\,f$ on
  $c(\sigma\times I, R)$, an equivalence if $f$ is; more generally a map
  $f : c(I,R) \to c(I',R)$ changing the fibre induces $\mathrm{blockMap}\,f$, and a family
  $T_{ij} : c(I,R)\to c(I',R)$ induces the rectangular block operator
  $\mathrm{blockOpMap}\,T$, with the expected composition rules against
  Proposition~\ref{prop:blockop}.
\end{definition}

\begin{theorem}[Perturbed unitriangular matrices are isometries]
  \label{thm:unitriangular}
  \uses{def:cspace}
  \lean{TateFredholm.IsUnitriangularPerturbation, TateFredholm.ofPerturbation,
        TateFredholm.norm_ofPerturbation, TateFredholm.surjective_ofPerturbation,
        TateFredholm.equivOfPerturbation}
  \leanok
  Let $M$ be an infinite matrix over a complete ultrametric field with entries of norm
  $\leq 1$, unit diagonal, entries strictly below the diagonal of norm $\leq q < 1$, and
  finitely supported columns.  Then $M$ defines an isometry of $c(\N,K)$ with dense range, hence
  a continuous linear equivalence.
\end{theorem}

\begin{proof}
  \leanok
  Norm-preservation is a largest-index argument: in $Mx$ the coordinate at the largest index
  where $x$ attains its norm receives the unit diagonal contribution and nothing else of that
  size, the sub-diagonal entries being smaller by the factor $q$.  Surjectivity is $q$-adic
  successive approximation.
\end{proof}

This is the criterion behind Colmez's proof of Amice's theorem, which Colmez states through the
residue field of a discretely valued base; the argument above needs only completeness and the
ultrametric inequality, which is what \S\ref{sec:amice} requires.

\begin{lemma}[Pairing of characteristic roots]
  \label{lem:charpoly-pairing}
  \lean{Matrix.det_mul_det_of_mul_eq_smul, Matrix.C_det_mul_charpolyRev_eq,
        Matrix.roots_charpolyRev, Matrix.roots_charpoly_of_mul_eq_smul}
  \leanok
  If $AB = c\cdot 1$ then $A$ is invertible with $B = c\,A^{-1}$, $\det A\det B = c^{\,n}$, and
  over an algebraically closed field
  $(\mathrm{charpoly}\,B).\mathrm{roots} = (\mathrm{charpoly}\,A).\mathrm{roots}$ mapped by
  $x \mapsto c/x$.
\end{lemma}

The root statement is phrased as an identity of multisets rather than as a sorted indexing
$\alpha_i = k+1-\alpha_{n-1-i}$: summing recovers the slope total, and counting recovers the
per-slope multiplicities, which are the two consumers in \S\ref{sec:stepone}.
